\documentclass[../../../include/open-logic-section]{subfiles}

\begin{document}

\olfileid{sth}{ord-arithmetic}{mult}

\olsection{د ترتيبي عددونو ضرب}

اوس د ترتيبي عددونو ضرب ته راځو، او د جمع په شان يې څېړو۔ فرض
کړئ چې $\alpha$ له~$\beta$ سره ضربول غواړو۔ د دې د تصور دپاره
يو مستطيل جال په پام کښې ونيسئ چې سور يې $\alpha$ او لوړوالے يې
$\beta$ وي؛ د $\alpha$ او $\beta$ حاصل ضرب دا دے چې د هر قطار
په اوږدو کښې لاړ شئ، بيا بل قطار ته ورشئ\ldots تر څو ټول جال
تېر کړئ۔ په بله وينا، د $\alpha$ او $\beta$ حاصل ضرب داسې جوړېږي
چې د $\beta$ \emph{هر} غړے د $\alpha$ په يوه کاپۍ بدل کړئ۔

د رسمي کولو دپاره د $\alpha$ او $\beta$ پر کارتيزي حاصل ضرب
هماغه برعکس قاموسي ترتيب کاروو:

\begin{defn}
د هر ډول ترتيبي عددونو $\alpha, \beta$ دپاره د دوی حاصل ضرب
$\alpha \ordtimes \beta = \ordtype{\alpha \times \beta, \rlexless}$ دے۔
\end{defn}
\noindent
بيا بايد پخلی وکړو چې دا تعريف په سمه توګه جوړ شوے دے:

\begin{lem}\ollabel{ordtimeslessiswo}
$\tuple{\alpha \times \beta, \rlexless}$ د هر ډول ترتيبي عددونو
$\alpha$ او $\beta$ دپاره ښه ترتيب دے۔
\end{lem}

\begin{proof}
عين د \olref[add]{ordsumlessiswo} په شان۔
\end{proof}
\noindent
دا هم ثابتول ګران نه دي چې ضرب لاندې چلند کوي:

\begin{lem}\ollabel{ordtimesrecursion}
د هر ډول ترتيبي عددونو $\alpha, \beta$ دپاره:
\begin{align*}
	\alpha \ordtimes 0 &= 0\\
	\alpha \ordtimes (\beta \ordplus 1) &=
		(\alpha \ordtimes \beta) \ordplus \alpha\\
	\alpha  \ordtimes \beta &=
		\supstrict_{\delta < \beta}(\alpha \ordtimes \delta) &&
		\text{چې $\beta$ حدي ترتيبي عدد وي}.
\end{align*}
\end{lem}

\begin{proof}
دا د تمرين دپاره پرېږدو۔
\end{proof}

په رښتيا، د جمع په شان مو ترتيبي ضرب هم د نېغ تعريف پر ځاے د
دغو بازګشتي معادلو له لارې تعريفولے شو۔ د جمع په څېر دلته هم
ځينې چلندونه اشنا دي:

\begin{lem}\ollabel{ordinalmultiplicationisnice}
که $\alpha, \beta, \gamma$ ترتيبي عددونه وي، نو:
\begin{enumerate}
	\item\ollabel{ordtimes1} که $\alpha \neq 0$ او $\beta < \gamma$
	وي، نو $\alpha \ordtimes \beta < \alpha \ordtimes \gamma$؛
	\item\ollabel{ordtimes2} که $\alpha \neq 0$ او $\alpha \ordtimes
	\beta = \alpha\ordtimes\gamma$ وي، نو $\beta = \gamma$؛
	\item\ollabel{ordtimes3} $\alpha \ordtimes (\beta \ordtimes
	\gamma) = (\alpha \ordtimes \beta) \ordtimes \gamma$؛
	\item\ollabel{ordtimes4} که $\alpha \leq \beta$ وي، نو $\alpha
	\ordtimes \gamma \leq \beta \ordtimes \gamma$؛
	\item\ollabel{ordtimes5} $\alpha \ordtimes (\beta \ordplus
	\gamma) = (\alpha \ordtimes \beta )\ordplus (\alpha\ordtimes
	\gamma)$۔
\end{enumerate}
\end{lem}

\begin{proof}
دا هم د تمرين دپاره پرېږدو۔
\end{proof}

نورې پايلې هم د خپلې خوښې تر بريده ثابتولے يا لوستلے شئ۔ خو د
\olref[ord-arithmetic][add]{ordsumnotcommute} له مخې لاندې خبره
بايد حيرانوونکې نه وي:

\begin{prop}
ترتيبي ضرب تبادلي نه دے: $2 \ordtimes \omega  =
\omega < \omega \ordtimes 2$۔
\end{prop}

\begin{proof}
$2 \ordtimes \omega = \supstrict_{n < \omega}(2\ordtimes  n) = \omega \in \supstrict_{n < \omega}(\omega \ordplus n) = \omega \ordplus \omega = \omega \ordtimes 2$۔
\end{proof}
\noindent
دلته هم بداهتي دليل ډېر روښانه دے۔ د $2
\ordtimes \omega$ د حساب دپاره هر طبيعي عدد په دوو څيزونو بدل کړئ۔
هماغه د ترتيب ډول ترلاسه کوئ که په طبيعي عددونو کښې د ټولو
«نيمو» عددونو ځايونه ورزيات کړئ؛ يعنې پر
$\Setabs{\nicefrac{n}{2}}{n \in \omega}$ طبيعي ترتيب وګورئ۔
په دې بڼه ئې د ترتيب ډول په څرګنده پخپله د $\omega$ په شان دے۔
خو د $\omega \ordtimes 2$ د حساب دپاره د $\omega$ دوې کاپۍ
يوه د بلې ورپسې ږدئ۔

\begin{prob}
\olref[sth][ord-arithmetic][mult]{ordtimeslessiswo}،
\olref[sth][ord-arithmetic][mult]{ordtimesrecursion}،
او \olref[sth][ord-arithmetic][mult]{ordinalmultiplicationisnice}
ثابت کړئ۔
\end{prob}

\end{document}
